\begin{trapbox}
	\textbf{\textcolor{CTrap}{易错一｜看见指数套对数，就直接“消掉”}}

	\textbf{错：}$4^{\log_2 3}=3$。
	\textbf{对：}$4=2^2$，所以
	\[
	4^{\log_2 3}=(2^{\log_2 3})^2=3^2=9.
	\]
	\textbf{原因：}外面的底数 $4$ 和对数底数 $2$ 不一样；\textbf{先改成同底，再还原}。

	\medskip
	\textbf{\textcolor{CTrap}{易错二｜忘记对数什么时候有意义}}

	\textbf{错：}认为 $2^{\log_2(x-1)}=x-1$ 对任意 $x$ 成立。
	\textbf{对：}必须先有 $x-1>0$，即 $x>1$，等式才有意义。
	同样，$2^{\log_2(-3)}$ 在实数范围内\textbf{没有意义}。

	\medskip
	\textbf{\textcolor{CTrap}{易错三｜忽略底数本身的要求}}

	对数的底数必须 $a>0$ 且 $a\neq 1$。
	例如 $\log_{-2}3$、$\log_1 3$ 在高中实数范围内都没有意义，不能套用恒等式。

	\medskip
	\textbf{\textcolor{CTrap}{易错四｜多了系数，也当成直接还原}}

	\textbf{错：}$2^{2\log_2 3}=3$。
	\textbf{对：}先把指数中的 $2$ 处理好：
	\[
	2^{2\log_2 3}=(2^{\log_2 3})^2=3^2=9.
	\]
	多出来的 $2$ 不能凭空消失。

	\medskip
	\textbf{\textcolor{CTrap}{对比记忆｜两个公式的条件别混}}
	\begin{center}
	\begin{tikzpicture}[>=stealth, font=\small]
		\node[draw=CTrap, rounded corners=3pt, text width=5.4cm, align=center, inner sep=5pt] at (0,0) {$a^{\log_a N}=N$\\$N>0$ 才能取对数};
		\node[draw=CTrap, rounded corners=3pt, text width=5.4cm, align=center, inner sep=5pt] at (6.05,0) {$\log_a(a^x)=x$\\$x$ 可取任意实数};
	\end{tikzpicture}

	\end{center}
	两边都要求 $a>0$、$a\neq 1$；第二式中 $a^x$ 本来就大于 $0$。
\end{trapbox}
